\chapter{Fundamentação Teórica}

A era contemporânea, frequentemente denominada como a ``era da pós-verdade'', representa uma mudança paradigmática onde fatos objetivos possuem menos influência na formação da opinião pública do que apelos a emoções e crenças pessoais \cite{umich_posttruth}. Neste capítulo, serão abordados os fundamentos técnicos tradicionais utilizados para a detecção de desinformação, suas limitações contemporâneas e a base matemática das Redes Neurais de Grafos (GNNs), que compõem o núcleo metodológico deste trabalho.

\section{Processamento de Linguagem Natural (NLP) e suas Limitações}

Historicamente, as primeiras tentativas de detecção automática de desinformação focaram nas propriedades intrínsecas do texto, partindo da premissa de que a linguagem enganosa possui marcadores sintáticos distintos das notícias legítimas \cite{pmc_slrfakenews, pmc_gnnsurveyfakenews}.

Modelos clássicos de aprendizado de máquina, utilizando técnicas como \textit{Bag of Words} (BOW) e TF-IDF para analisar a frequência de termos, aliados a algoritmos como \textit{Random Forest} e \textit{Support Vector Machines} (SVM), alcançam precisão entre 80\% e 83\% em conjuntos de dados padronizados \cite{scitepress_nlpfakenews}. A evolução para modelos de linguagem pré-treinados baseados na arquitetura \textit{Transformer}, como o BERT (\textit{Bidirectional Encoder Representations from Transformers}), permitiu a compreensão do contexto bidirecional das palavras, elevando as taxas de acerto \cite{mdpi_cnn_llm_nlp, pmc_slrfakenews}.

Contudo, a dependência exclusiva do conteúdo textual torna esses modelos vulneráveis, especialmente com o advento de \textit{deepfakes} textuais capazes de mimetizar o estilo jornalístico \cite{pmc_multilangs_covid, mdpi_cnn_llm_nlp}. A literatura aponta para um ``abismo de generalização'' com três frentes críticas: a mudança de idioma, a assimetria de recursos e a variação entre registros formais e informais \cite{pmc_multimodalfakenews, stanford_nlpmisinformation}.

\subsection{O Abismo Linguístico e o Contexto Brasileiro}

A detecção focada apenas em NLP sofre perdas drásticas de precisão ao ser transposta para outros idiomas sem dados de treinamento nativos \cite{crosslanguagefakenews}. Estudos demonstram fortes assimetrias de desempenho baseadas na complexidade da língua de origem durante o treinamento translinguístico \cite{crosslanguagefakenews, cogent_nlpmisinformation}.

Além disso, há um hiato estilístico severo. Enquanto notícias legítimas operam em linguagem formal, a desinformação em redes sociais adota uma linguagem coloquial, repleta de gírias e regionalismos \cite{cogent_nlpmisinformation}. No Brasil, pesquisas apontam que modelos treinados em bases de agências de checagem frequentemente falham em campo, pois o texto metodológico do checador não reflete a estrutura orgânica da \textit{fake news} real que circula em mensageiros e redes sociais \cite{ufscar_iabrfakenews, ufc_fakewhatsapp}.

\section{Teoria de Grafos e a Geometria da Disseminação}

Diante das limitações do processamento estritamente textual, a fronteira de pesquisa deslocou-se para a estrutura de disseminação. A premissa central é que a verdade e a mentira possuem ``formas'' geométricas diferentes na maneira como se propagam \cite{pmc_gnnsurveyfakenews, monti2019fakenews}.

Estudos empíricos confirmam que notícias falsas penetram em cascata por muitos níveis de compartilhamento e atingem milhares de pessoas de forma consideravelmente mais rápida que notícias reais \cite{pnasnexus_misinformation, researchgate_temporalgraphs}. Esse fenômeno é impulsionado por ``câmaras de eco'': comunidades ideologicamente homogêneas, com altos coeficientes de agrupamento (\textit{clustering coefficients}), onde a informação falsa circula e é reforçada continuamente \cite{pnasnexus_misinformation}.

Matematicamente, a rede de propagação é modelada como um grafo $G = (V, E)$ \cite{pmc_gnnsurveyfakenews}. Neste contexto, $V$ representa o conjunto de nós, que podem ser os usuários ou as publicações, e $E$ representa o conjunto de arestas, que mapeiam as interações como \textit{follows}, \textit{likes} e \textit{reposts} \cite{pmc_gnnsurveyfakenews, arxiv_topologicalfeatures}.

\section{Redes Neurais de Grafos (GNN)}

As Redes Neurais de Grafos operam diretamente sobre a estrutura topológica, integrando o conteúdo textual (vetorizado nos nós) com a estrutura social \cite{pmc_gnnsurveyfakenews, monti2019fakenews}.

\subsection{Graph Convolutional Networks (GCN)}

O modelo GCN, proposto por \citeonline{kipf2017gcn}, utiliza uma aproximação de convoluções espectrais em grafos para aprender as representações espaciais dos nós. A regra fundamental de propagação camada a camada em uma GCN é expressa matematicamente pela equação:

\begin{equation}
H^{(l+1)} = \sigma \left( \tilde{D}^{-\frac{1}{2}} \tilde{A} \tilde{D}^{-\frac{1}{2}} H^{(l)} W^{(l)} \right)
\end{equation}

Onde:
\begin{itemize}
    \item $H^{(l)}$ é a matriz de representação (estado oculto) da camada atual \cite{kipf2017gcn};
    \item $\tilde{A} = A + I_N$ é a matriz de adjacência do grafo com a adição de auto-conexões, garantindo que o nó considere sua própria informação \cite{kipf2017gcn};
    \item $\tilde{D}$ é a matriz de graus utilizada para a normalização simétrica \cite{kipf2017gcn};
    \item $W^{(l)}$ é a matriz de pesos treináveis da respectiva camada \cite{kipf2017gcn};
    \item $\sigma$ denota a função de ativação não linear (como ReLU) \cite{kipf2017gcn}.
\end{itemize}

Essa arquitetura apresenta alta eficácia em cenários semi-supervisionados, sendo ideal para redes sociais onde apenas uma fração das publicações possui rótulos de veracidade confirmados \cite{pmc_pulpfakenews}.

\subsection{Graph Attention Networks (GAT)}

Como extensão evolutiva, as Redes de Atenção em Grafos (GAT), propostas por \citeonline{velickovic2018gat}, introduzem mecanismos de atenção para atribuir pesos dinâmicos às conexões. Diferente da GCN tradicional, a GAT reflete a realidade de que nem todos os usuários exercem a mesma influência na rede. O coeficiente de atenção $e_{ij}$ entre o nó $i$ e o nó $j$ é dado por:

\begin{equation}
e_{ij} = \text{LeakyReLU} \left( a^T [Wh_i \parallel Wh_j] \right)
\end{equation}

Onde $\parallel$ representa a concatenação dos vetores transformados dos nós. Essa flexibilidade permite que o modelo foque automaticamente em ``nós hubs'' ou padrões de interação suspeitos, garantindo maior resiliência a variações topológicas da plataforma analisada \cite{velickovic2018gat, cogent_nlpmisinformation}.
